\chapter{Literatura i podsumowanie}
\label{ch:kt-literatura-podsumowanie}

Wcześniejsze wydanie skryptu zawierało osobny rozdział ze szczegółowymi rachunkami. W bieżącej wersji obliczenia zostały wplecione bezpośrednio w odpowiednie części wykładu: grupa trójlistnika pojawia się przy prezentacji Wirtingera, rachunki z macierzy Seiferta przy wielomianie Alexandera i sygnaturze, obliczenie wielomianu Jonesa w sekcji o nawiasie Kauffmana, przykład przestrzeni soczewkowych przy chirurgii Dehna, a obliczenia formy przecięcia przy uchwytach w wymiarze cztery. Dzięki temu aparat obliczeniowy występuje dokładnie tam, gdzie jest używany.

\section{Literatura podstawowa i zaawansowana}
\label{sec:kt-literatura}

\begin{enumerate}[label={[\arabic*]},leftmargin=*,itemsep=3pt]
\item D. Rolfsen, \emph{Knots and Links}, AMS Chelsea Publishing.
Wprowadzenie klasyczne: diagramy, grupy, wielomiany i sploty.
\item C. Livingston, \emph{Knot Theory}, Carus Mathematical Monographs.
Wykład z naciskiem na rachunki i powierzchnie Seiferta.
\item W. B. R. Lickorish, \emph{An Introduction to Knot Theory}, Springer.
Wielomiany Jonesa, skein i chirurgia Dehna.
\item G. Burde, H. Zieschang, M. Heusener, \emph{Knots}, De Gruyter.
Grupy, rozwarstwienia i klasyczna topologia uzupełnień.
\item J. Hempel, \emph{3-Manifolds}, AMS Chelsea.
Tło dla powierzchni nieściśliwych i topologii trójrozmaitości.
\item W. P. Thurston, \emph{The Geometry and Topology of Three-Manifolds}.
Geometria hiperboliczna i wypełnienia Dehna.
\item R. E. Gompf, A. I. Stipsicz, \emph{4-Manifolds and Kirby Calculus}, AMS.
Podstawowa monografia o uchwytach, chirurgii i formach przecięcia.
\item M. H. Freedman, F. Quinn, \emph{Topology of 4-Manifolds}, Princeton.
Topologiczna kategoria wymiaru cztery i zanurzanie dysków.
\item P. Ozsváth, Z. Szabó, \emph{Holomorphic disks and knot invariants}.
Konstrukcja homologii Floera węzłów.
\item M. Khovanov, \emph{A categorification of the Jones polynomial}.
Pierwotna konstrukcja homologii Khovanova.
\item J. Rasmussen, \emph{Khovanov homology and the slice genus}.
Niezmiennik $s$ i rozwiązanie problemu Milnora metodami homologii.
\item P. Kronheimer, T. Mrowka, \emph{Khovanov homology is an unknot-detector}.
Detekcja węzła trywialnego przy użyciu teorii instantonów.
\item P. Orlik, \emph{Seifert Manifolds}, Lecture Notes in Mathematics 291.
Klasyczne rozwarstwienia okręgowe, grupy i ich klasyfikacja.
\item W. D. Neumann, \emph{A calculus for plumbing applied to the topology of complex surface singularities}.
Rachunek ważonych grafów, ich operacji i brzegów.
\item M. Hedden, \emph{Knot Floer homology of Whitehead doubles}.
Obliczenia $\tau$ i rozróżnianie slicowości gładkiej i topologicznej.
\item R. Fintushel, R. J. Stern, \emph{Rational Blowdowns of Smooth
$4$-Manifolds}, J. Differential Geom. 46 (1997), 181--235.
Wprowadzenie wymiernej kontrakcji i formuły dla niezmienników gładkich.
\item A. Casson, J. Harer, \emph{Some homology lens spaces which bound
rational homology balls}, Pacific J. Math. 96 (1981), 23--36.
Konstrukcje wymiernych kul homologii.
\item J. Park, \emph{Simply connected symplectic $4$-manifolds with $b_2^+=1$ and $c_1^2=2$}, Invent. Math. 159 (2005), 657--667. Konstrukcja $C_7$ w $E(1)\#4\overline{\mathbb{CP}}^{\,2}$ oraz argument prostej spójności dopełnienia. \url{https://doi.org/10.1007/s00222-004-0404-1}
\item J. Park, A. I. Stipsicz, Z. Szabó, \emph{Exotic smooth structures
on $\mathbb{CP}^2\#5\overline{\mathbb{CP}}^{\,2}$},
Math. Res. Lett. 12 (2005), 701--712.
Przykłady uogólnionych wymiernych kontrakcji.
\item J. Park, \emph{Seiberg--Witten invariants of generalised rational
blow-downs}, Bull. Austral. Math. Soc. 56 (1997), 363--384.
Zależności niezmienników od warunków przedłużenia klas.
\item R. Fintushel, R. J. Stern, \emph{Six Lectures on Four-Manifolds},
IAS/Park City Math. Ser. 15 (2009), 265--315.
Wprowadzenie do niezmienników Seiberga--Wittena i racjonalnych kontrakcji.
\item J. W. Morgan, \emph{The Seiberg--Witten Equations and Applications
to the Topology of Smooth Four-Manifolds}, Princeton University Press.
Analityczne podstawy teorii cechowania i konstrukcja niezmienników.
\item R. Fintushel, R. J. Stern, \emph{Knots, Links, and
$4$-Manifolds}, Invent. Math. 134 (1998), 363--400.
Oryginalne twierdzenie o chirurgii wzdłuż węzła, formuła
Seiberga--Wittena i konstrukcje egzotycznych rozmaitości K3.
\url{https://arxiv.org/abs/dg-ga/9612014}
\item R. Fintushel, R. J. Stern, \emph{Constructions of Smooth
$4$-Manifolds}, Documenta Math. Extra Volume ICM 1998.
Dodatkowe objaśnienia formuły chirurgii wzdłuż węzła.
\item T.-J. Li, A.-K. Liu, \emph{General wall crossing formula},
Math. Res. Lett. 2 (1995), 797--810.
Klasyczna formuła zmiany komory dla $b_2^+=1$.
\item L. I. Nicolaescu, \emph{Notes on Seiberg--Witten Theory},
Graduate Studies in Mathematics, AMS. Oszacowania,
przestrzenie moduli, twierdzenia o zwartości i zmiana komory.
\end{enumerate}

\section*{Podsumowanie zależności}
Najważniejsze przejścia, które warto zachować przy ponownym
wykorzystaniu skryptu w rozdziałach o wymiarze cztery, są następujące:
\begin{enumerate}[label=(\roman*)]
\item diagram węzła $\longrightarrow$ prezentacja Wirtingera,
powierzchnia Seiferta, macierz Seiferta i wielomian Alexandera;
\item diagram zorientowany $\longrightarrow$ nawias Kauffmana,
wielomian Jonesa oraz kompleks Khovanova;
\item uzupełnienie węzła $\longrightarrow$ torus peryferyjny,
nachylenia, chirurgia Dehna i nowe zamknięte trójrozmaitości;
\item ramowany splot $\longrightarrow$ uchwyty indeksu dwa,
forma przecięcia i rachunek Kirby'ego w wymiarze cztery;
\item dysk w $B^4$ $\longrightarrow$ gładka i topologiczna
konkordancja, twierdzenie Freedmana oraz niezmienniki
$\tau$ i $s$ rozróżniające kategorie TOP i DIFF;
\item włókno torusowe $T\subset E(2)$ oraz węzeł $K\subset S^3$
$\longrightarrow$ chirurgia Fintushela--Sterna i formuła
$\mathcal{SW}_{E(2)_K}=\Delta_K(\exp(2[T_K]))$,
prowadzące do egzotycznych struktur na K3.
\end{enumerate}

% Koniec tresci w wersji zintegrowanej. Przy wstawieniu do wiekszej ksiazki
% rozdzialy sa automatycznie numerowane zgodnie z jej struktura.


